\section{Related Work}

\subsection{ADMM and sharing for demand response}
ADMM \cite{boyd2011admm} decomposes a coupled optimization into local subproblems
linked by a consensus or sharing constraint, and has become a standard tool for
distributed energy coordination. The sharing form (Boyd et al., Section~7.3) is
particularly suited to demand flattening, where many agents contribute to a
single shared resource. Asynchronous and energy-grid-specific variants extend it
to microgrids and optimal power flow
\cite{UllahPark2019,Li8409328}. These works establish that ADMM converges to a
flattened aggregate; they do not, in general, check that the aggregate is
feasible on the underlying network.

\subsection{Multi-agent resilience to non-responsiveness}
A separate line of work studies how multi-agent coordination degrades when
agents lose communication. Brown et al. show that even a single lost link can
drive game-theoretic multi-agent systems to arbitrarily poor states
\cite{Brown2017}, and consensus-control surveys catalog fault-tolerant
mechanisms \cite{GaoHe2022}. In the distributed-optimization setting,
communication delay and loss have been shown to slow or destabilize ADMM-based
optimal power flow \cite{Guo2017}. This literature characterizes \emph{when}
coordination breaks but rarely couples the failure model to a concrete recovery
mechanism whose error is then propagated into feeder physics.

\subsection{Imputation of missing telemetry}
When agents fall silent, the coordinator must estimate their contribution.
Machine-learning imputation of missing time-series telemetry, and its downstream
effect on forecasting, is an active area \cite{Saad2020}, with cooperative
network estimators able to fill missing data across agents
\cite{Gholami2016}. Gradient-boosted trees such as LightGBM \cite{ke2017lightgbm}
are strong tabular regressors and are well matched to imputing a single agent's
heating from weather and calendar features.

\subsection{Distribution power flow}
Physical validation requires a feeder model and a solver. The IEEE 33-bus radial
test feeder \cite{baran1989network} is the canonical distribution benchmark, and
pandapower \cite{thurner2018pandapower} provides convenient AC Newton-Raphson
power flow over it. These tools are mature; what is missing is their systematic
use as the \emph{acceptance test} for a coordination scheme.

\subsection{Communication in demand-side management}
The reliability of the communication layer is increasingly recognized as a
first-order constraint on demand-side management rather than an implementation
detail. Surveys of information and communication technologies for modern
demand-side management catalog the latency, packet-loss, and coverage limitations
of the metering and signalling infrastructure that coordination schemes assume
\cite{Said2023}, and identify intermittent connectivity behind the meter as a
recurring obstacle to scalable demand response. This motivates our explicit
non-responsive-fraction model: rather than assuming perfect reporting, we treat
silent agents as the expected operating condition and ask how far coordination
can be pushed before the missing telemetry --- even with a learned imputer in the
loop --- can no longer protect the feeder.

\subsection{The delta}
Prior work treats coordination quality and network feasibility as separate
concerns: ADMM papers report aggregate metrics, network papers assume a given
load. We connect them. By feeding the coordinated (and imputation-perturbed)
aggregate into an AC power flow and reading back thermal and voltage violations,
we measure coordination by the metric an operator cares about: does the feeder
stay within limits? To our knowledge this closed loop --- ADMM, imputation, and
power-flow validation under a swept communication-failure fraction --- has not
been reported for cold-climate electric heating.
